% Generated draft: requires mathematical and source review.
\documentclass[11pt]{article}
\usepackage[margin=25mm]{geometry}
\usepackage{fontspec}
\setmainfont{Latin Modern Roman}
\usepackage{amsmath,amssymb,mathtools,unicode-math}
\setmathfont{Latin Modern Math}
\usepackage{enumitem,needspace,xurl,hyperref}
\setlength{\parindent}{0pt}
\setlength{\parskip}{5pt}
\setlength{\emergencystretch}{3em}
\newcommand{\N}{\mathbb{N}}
\newcommand{\Z}{\mathbb{Z}}
\newcommand{\Q}{\mathbb{Q}}
\newcommand{\R}{\mathbb{R}}
\newcommand{\C}{\mathbb{C}}
\newcommand{\F}{\mathbb{F}}
\newcommand{\A}{\mathbb{A}}
\newcommand{\PP}{\mathbb{P}}
\newcommand{\E}{\mathbb{E}}
\newcommand{\eps}{\varepsilon}
\newcommand{\dd}{\,\mathrm{d}}
\begin{document}
\section*{UnsolvedMath ID: 4600051}
\textbf{Source problem code:} catalogue entry
\subsection{Definitions and mathematical statement}
A subshift in the usual two-sided finite-alphabet setting is a closed shift-invariant subset \(X\subseteq\mathcal A^{\Z}\), where \(\mathcal A\) is finite and \((\sigma x)_j=x_{j+1}\). Write \((M_X,I_X)\) for its canonical matrix system in the intended source convention. The record does not define that construction or the associated groups \(K_0(M_X,I_X)\) and \(K_1(M_X,I_X)\).

The requested range classification is to characterize exactly the ordered pairs of abelian-group isomorphism classes
\[
\mathcal R=\{([K_0(M_X,I_X)],[K_1(M_X,I_X)]):
 X\text{ is an admissible subshift}\}.
\]
Equivalently, give necessary and sufficient conditions on abelian groups \(G_0,G_1\) for
\[
\exists X\quad K_0(M_X,I_X)\cong G_0
\quad\text{and}\quad K_1(M_X,I_X)\cong G_1.
\]
The two groups must be realized by the same canonical system. The question concerns their underlying abelian groups; no order cone, distinguished unit, or additional invariant is included in the supplied statement. No restriction to finite-type shifts or finitely generated groups is stated. Realization by an arbitrary noncanonical matrix system would not alone answer the question.

The canonical-system and K-group definitions, and the precise admissible subshift class, must be supplied before this range problem becomes fully determined. In particular, a finite-matrix kernel/cokernel formula or a chosen associated operator algebra is not silently substituted for these source-specific groups. The supplied research summary reports the full range problem open through 2026 and relates it to dimension-group and operator-algebraic classification, without supplying a characterization.
\subsection{Short English statement}
Which pairs of abelian groups can be realized simultaneously as the two K-groups of one subshift's canonical matrix system? The canonical construction and admissible class must be fixed to specify the range.
\subsection{Sources}
\begin{itemize}
\item Local UnsolvedMath registry, permanent ID \(4600051\), exact source code \texttt{catalogue entry}, ``K-groups of canonical matrix systems.''
\item Boyle, \emph{Open Problems in Symbolic Dynamics} (2008), Problem/Question/Conjecture 33.2; \url{https://www.math.umd.edu/~mboyle/papers/openfinalsub3nov2007.pdf}. Recorded access 2026-07-29; pdf-text extraction.
\item Supplied OPEN-TRIAGE research summary mentions Krieger and Boyle--Handelman in relation to dimension groups and K-theoretic classification. It supplies no complete publication references or definitions of the specific canonical-system K-groups.
\end{itemize}
\end{document}
